\chapter{Introduction}\label{ch:intro}

\begin{comment}
This chapter: 
DM and CASPEr need NMR.
->what type of NMR do we need?
low-noise, any field (ultra-low?)
\end{comment}


\section{Dark matter research with nuclear spin}
% Vera Rubin, Kent Ford, and Ken Freeman's work in the 1960s and 1970s[36] provided further strong evidence, also using Galaxy rotation curves

Various astrophysical observations implied that with ``visible'' mass and known theories on gravity, some phenomena cannot be explained rationally. Take Galaxy rotation curve as an example. The Galaxy does not fall apart during its rotation due to the gravitational effect. However, if we do not work in the direction of modified Newtonian dynamic (MOND) or other new ideas, we have to admit that there has been ``invisiable'' massive matter in the universe. One of such example is the rotation curve of M33 halo\cite{Corbelli_M33}. Therefore the concept of ``dark matter" (DM) was invented. 

The properties of dark matter is an open question. Although total amount of dark matter is significant in comparison to known matters (see Fig.\,\ref{fig:matter_mass_ratio}), the average density of it in the universe is still too low to be detected directly through gravitational effect. Thus we expect other forms of interaction. 
\begin{figure}[htbp]
    \centering
    \resizebox{0.8\linewidth}{!}{
    \begin{tikzpicture}
        \pie[text=pin, explode={0.0, 0.0, .0 , 0.0}, hide number]{73/{Dark energy [73\%]}, 23/{Dark matter [23\%]}, 3.6/{Intergalactic gas [3.6\%]}, 0.4/{Stars, etc [0.4\%]}}
    \end{tikzpicture}}
    \caption{Estimated mass distribution in the universe\cite{Ackermandmater}. }
    \label{fig:matter_mass_ratio}
\end{figure}

Axion, as a hypothetical elementary particle, was first invented in 1977 as a solution to the strong $CP$ problem in quantum chromodynamics (QCD). %Sometimes it is also named as QCD axion. 
It was found to be a possible candidate of dark matter since it could acquire mass from QCD when some global symmetry is broken at high energy scale $f_a$. The concept of axion-like particle (ALP) was later created to classify those which can acquire mass from mechanisms other than QCD, hence the possible mass range of ALPs is expected to be much larger than QCD axion. To simplify the text, we do not always differentiate these two concepts and use ``axion'' to refer to both of them. 

There are multiple hypothesises on the interaction between axion and other particles\cite{GrahamNewobservables_PhysRevD.88.035023}, including axion-photon, axion-gluon and axion-fermion. Cosmic axion spin precession experiment (CASPEr) was proposed in 2014\cite{budkergrahamproposal} to explore the last two of potential interactions. 
The project has two branches: CASPEr-electric and CASPEr-gradient. The first searches for axion-gluon coupling and the other aims at axion-fermion coupling. Besides, CASPEr-gradient was planned to have low-field and high-field phases, using \SI{0.1}{\tesla} and \SI{14.1}{\tesla} magnet respectively. In this thesis, we mainly talk about low-field phase. 

\section{Axion wind and pseudo-magnetic field}\label{sec:Axion_wind}

\subsection{Brief description of axion field}\label{sec:Axion_wind:description}

With interactions of axion field and initial displacement from the minimum potential, oscillation is naturally expected for the axion field. Compton frequency $\omega_a$ can be computed as:
\begin{equation}\label{eq:axion_omega_ma}
    \omega_a = \dfrac{m_a c^2}{\hbar}\,. 
\end{equation}

Here we describe axion particle as a classical oscillator with expression
\begin{equation}
    a = a_0 \cos(\omega_a t - \mathbf{k}\cdot \mathbf{r} + \phi)\label{eq:axion_field}\,,
\end{equation}
where $a_0$ is the amplitude of the field, $\mathbf{k} = m_a \mathbf{v}/\hbar$ is the wave vector of the axion particle, $r$ is the displacement vector, and $\phi$ is a random phase on $[0,\, 2\pi)$. The amplitude $a_0$ can be estimated from the local dark matter energy density $\rho_{DM} \approx \SI{0.4}{\GeV \cm^{-3}}$\cite{BERGSTROM1998137, Jungman1992SupersymmetricDM, Sofue_RotationCurves} in this way:
% 
\begin{align}
    \dfrac{c m_a^2 a_0^2}{2\hbar^3} \approx \rho_{DM}\,,\label{eq:a0estimation} 
\end{align}
where the left side is the potential energy density generated by the interactions of axion field. Note that the $a$ possesses the unit of energy. 

Nucleon is predicted to be able to interact with axion by emitting it or exchanging it with other nucleons\cite{GrahamNewobservables_PhysRevD.88.035023}. Nucleon Hamiltonian with axion field background can be written as
\begin{align}
    H_a = \mathrm{g_{aNN}} {\boldsymbol \nabla} a \cdot \mathbf{I} 
    = \hbar \mathrm{g_{aNN}} {\boldsymbol \nabla} a \cdot \hat{\boldsymbol \sigma}_{N}\,,\label{eq:axion_nucleonH}
\end{align}
where $\mathrm{g_{aNN}}$ is the coupling constant in the unit of \SI{}{\GeV^{-1}} and $\hat{\boldsymbol \sigma}_{N}$ is the dimensionless nucleon spin operator. 

With the presence of magnetic field, the nucleon Hamiltonian acquires:
\begin{align}
    H_{\mathrm{NMR}} = \gamma \mathbf{B} \cdot \mathbf{I}
    = \hbar \gamma \mathbf{B} \cdot \hat{\boldsymbol \sigma}_{N}\,,\label{eq:magfield_nucleonH}
\end{align}
where $\gamma$ is the nuclear gyromagnetic ratio and $\mathbf{B}$ is the magnetic field. From the comparison between Eq.\,\eqref{eq:axion_nucleonH} and \eqref{eq:magfield_nucleonH}, we can find the analogy from axion spatial gradient to magnetic field. This gives us the hint that we can regard axion field as ``pseudo-magnetic field'' and study it with mature tools for magnetic field measurement. We can rewrite Eq.\,\eqref{eq:axion_nucleonH} into
\begin{equation}\label{eq:pseudo-magfield_nucleonH}
    H_a =\hbar \gamma \mathbf{B}_{a}(t) \cdot \hat{\boldsymbol \sigma}_{N}\,,
\end{equation}
where $\mathbf{B}_{a}(t)$ is the time-dependent pseudo-magnetic field. The explicit form for $\mathbf{B}_{a}(t)$ can be inferred from Eq.\,\eqref{eq:axion_field} and \eqref{eq:axion_nucleonH} (neglecting the minus symbol):
\begin{equation}
    \mathbf{B}_{a}(t) = \mathrm{g_{aNN}} {\boldsymbol \nabla}a/ \gamma = \mathrm{g_{aNN}} m_a a_0 \mathbf{v}_a \sin(\omega_a t - \mathbf{k}\cdot \mathbf{r} + \phi) / \gamma\,, \label{eq:Bat_a0}
\end{equation}
With Eq.\,\eqref{eq:a0estimation}, the amplitude of $\mathbf{B}_{a}(t)$ can be estimated as
\begin{align}
    \mathrm{B}_{a} &\approx \mathrm{g_{aNN}} \sqrt{2 \hbar c \rho_{DM}} v_a / \gamma  
    \approx  \mathrm{g_{aNN}} \mathrm{g_{DM}}  / \gamma 
    \,.\label{eq:Ba_gaNN_rhoDM}
    % &\approx 120 [\SI{}{\GeV .\Hz}] \times \dfrac{\mathrm{g_{aNN}}}{\gamma} \,,
\end{align}
Here $\mathrm{g_{DM}} \equiv \sqrt{2 \hbar c \rho_{DM}} v_a \approx \SI{4}{\GeV.\Hz}$. Assuming there is no significant variation in the local dark matter density, the unknown coupling constant determines the strength of pseudo-magnetic field. 
% 2\pi\times\SI{0.6}{\GeV.\Hz}
\subsection{Stochastic nature}\label{sec:Axion_wind:Stochastic_nature}

By far, we have obtained a preliminary understanding of the axion field. 
% The axion field is a collective effect of axion ensemble. 
To better depict axion ensemble, we introduce such an expression for axion field:
\begin{equation}
    a = \dfrac{a_0}{\sqrt{N_a}}\sum_{n=1}^{N_a}  \cos(\omega_n t - \mathbf{k}_n \cdot \mathbf{r} + \phi_n)\,.\label{eq:axion_ensemble}
\end{equation}
Here $N_a$ is the number of axion particles, $\omega_n$ is the frequency of n-th particle, $\mathbf{k}_n$ is the wave vector, and $\phi_n$ is a random phase on $[0,\, 2\pi)$. 
% Two crucial terms in, though, have not yet been taken into account. 
% We can further explore the properties of $\omega_n$ and $\phi_n$ here. 

The gradient of axion field can be inferred as:
\begin{equation}
    {\boldsymbol \nabla}a =  \dfrac{a_0 m_a }{\sqrt{N_a} \hbar}\sum_{n=1}^{N_a} \mathbf{v}_n \sin(\omega_n t - \mathbf{k}_n\cdot \mathbf{r} + \phi_n)\,. \label{eq:axion_ensemble_gradient}
\end{equation}
We can go over some of the parameters of the expression here which are involved in the stochastic nature of axion. 
% \begin{enumerate}
    
\textbf{Velocity}

According to the standard halo model (SHM) for the DM halo of our Galaxy\cite{Schumann_2019, Evans_PhysRevD.99.023012}, the velocity distribution of DM particles in the galactic rest frame can be characterized by the Maxwell-Boltzmann distribution along with a Heaviside function:
\begin{equation}
    f(\mathbf{v})d^3\mathbf{v} = \dfrac{1}{\pi^{3/2} v_0^3 N_{esc}} e^{- \mathbf{v}^2 / v_0^2} \Theta(v_{esc} - |\mathbf{v}|)d^3\mathbf{v}\,,
\end{equation}
where the value of $v_0 \approx \SI{220}{\km / \second}$ is usually taken as the circular rotation speed of the Galaxy at the Solar position, and $v_{esc} \approx \SI{544}{\km / \second}$ is the Galaxy escape speed. Heaviside function $\Theta$ indicates that particles moving faster than escape speed are not bounded by the gravitational potential of the Galaxy. In order to renormalize the probability distribution function, $N_{esc}$ is introduced:
\begin{equation}
    N_{esc} = \mathrm{erf}\left(\dfrac{v_{esc}}{v_0}\right) - \dfrac{2}{\sqrt{\pi}} \dfrac{v_{esc}}{v_0}e^{- v_{esc}^2 / v_0^2}\,.
\end{equation}
It has to be mentioned that, velocity can be influenced by massive items through gravity. For example, Sun can focus incoming DM particles and modulate the density and velocity of that. 
    
\textbf{Frequency}

$\omega_n$ varies from axion Compton frequency $\omega_a$ because of Doppler effect. The frequency a detector on Earth sees is
\begin{equation}\label{eq:axion_velocity_Doppler}
    \omega_n = \omega_a \left(1 + \dfrac{v_n^2}{c^2}\right)\,. 
\end{equation}
Take $v_n \approx 10^{-3}c$ we have $v_n^2/c^2 \approx 10^{-6}$. 
From Eq.\,\eqref{eq:axion_velocity_Doppler} we know that the spread of velocities turns into the spread in frequency.  
    
\textbf{Detector position}
    
The position vector $\mathbf{r}$ is the displacement of the detector. In the case of $n_d$ detectors located at $\mathbf{r}_{i}$ ($i = 1,\,2\,, \ldots n_d$), there are phase differences of $\mathbf{k} \cdot \Delta \mathbf{r}$ between detectors. 

We have to answer such a question: what is the threshold for the distance between detectors to spatially distinguish them in terms of the phase difference of axion field? This question is vital not only for multiple-detector scenario but also for phase-sensitive detector which is not small enough to be represented by one point. Here we adopt $\Delta r$ to stand for the dimension of detector. The condition for insignificant phase difference is
\begin{equation}
    \mathbf{k} \cdot \Delta \mathbf{r} \ll 2\pi\,.\label{eq:k_Deltar}
\end{equation}
% The direction of the velocity and wave vector is modulated by the motion of Sun and Earth. 
The angle between $\mathbf{k}$ and $\Delta \mathbf{r}$ can be an arbitrary value. Here we can consider $\mathbf{k} \cdot \Delta \mathbf{r} = |\mathbf{k}| |\mathbf{r}|$. If we plug in $\mathbf{k} = m_a \mathbf{v}/\hbar$, $m_a/\hbar = \omega/c^2$ and $|\mathbf{v}|\sim10^{-3}c$, Eq.\,\eqref{eq:k_Deltar} turns into
\begin{equation}
    |\Delta \mathbf{r}| \ll 10^{3} \dfrac{c}{\nu_a}\,. \label{eq:dimension_limit_condition}
\end{equation}
Some values for $10^{3} c/\nu_a$ are listed in Tab.\,\ref{tab:dimension_limit}. We can take \SI{1}{\percent} of this factor as the threshold for neglecting phase difference $\mathbf{k} \cdot \Delta \mathbf{r}$. The Earth diameter is approximately \SI{1.28e7}{\metre}. For multi-detector experiment on Earth, the phase difference starts to play a role from $\sim$\SI{100}{\Hz}. On the other hand, for a common laboratory-size experiment with one single detector ($|\Delta \mathbf{r}| < \SI{3}{\metre}$), it does not matter until $\sim$\SI{1}{\GHz}. 
% CASPEr-Gradient adopts sample holder whose dimension is $\sim$\SI{1}{\centi\metre}. Hence we can use

\begin{table}[htbp]
\centering
\caption{Threshold for neglecting axion field phase difference between spatially distributed detectors}
\begin{tabular}{cccccc}
	\hline
	$\nu_a$ & \SI{1}{\mHz} & \SI{1}{\Hz} & \SI{1}{\kHz} & \SI{1}{\MHz} & \SI{1}{\GHz} \\
	\midrule
	$10^{3} \dfrac{c}{\nu_a}$ / \SI{}{\metre} & \SI{3e14}{} & \SI{3e11}{} & \SI{3e8}{} & \SI{3e5}{} & \SI{3e2}{} \\
	\midrule
	$10^{1} \dfrac{c}{\nu_a}$ / \SI{}{\metre} & \SI{3e12}{} & \SI{3e9}{} & \SI{3e6}{} & \SI{3e3}{} & \SI{3}{} \\
	\bottomrule
\end{tabular}
\label{tab:dimension_limit}
\footnotesize{}
\end{table}

\textbf{Phase}

$\phi_n$ does not remain as a constant value all the time, for which reason the phase of ${\boldsymbol \nabla}a$ (see Eq.\,\ref{eq:axion_field}), $\phi$, is not deterministic. We characterize it by coherence time $\tau_a$. 
The process of decoherence is gradual, slow while destructive. 
Within any short period ($\ll \tau_a$), we should not observe significant change in the phase of ${\boldsymbol \nabla}a$. In contrast, when measuring the ensemble at interval of or larger than $\tau_a$, we can find by the end of this period, $\phi$ is different from the beginning. 
Such decoherence can be regarded as the result of accumulating phase jumps of single particles over time. 
``Thermal light source'' or equivalently ``chaotic light source''\cite{Loudon2000TheQT} can be a good model to describe such decoherence process. 

We can estimate $\tau_a$ by its de Broglie wavelength $\lambda_{dB}$ and the speed $v_a$:
\begin{equation}
    \tau_a \approx \dfrac{\lambda_{dB}}{v_a} \approx \dfrac{2\pi \hbar}{m_a v_a^2}\,. \label{eq:taua_ma_va}
\end{equation}
Furthermore, a $Q$ factor can be derived from Eqs.\,\eqref{eq:axion_omega_ma} and \eqref{eq:taua_ma_va} to characterize the relationship between frequency and coherence time:
\begin{equation}
    Q = \dfrac{\omega_a \tau_a}{2 \pi} = \nu_a \tau_a \approx \left(\dfrac{c}{v_a}\right)^2 \approx 10^6\,. \label{eq:axion_Q_factor}
\end{equation}
One way to interpret the $Q$ factor is the relationship between axion decoherence rate $1/\tau_a$ and frequency $\omega$:
\begin{equation}
    \dfrac{1}{\tau_a} = 10^{-6} \nu_a\,.\label{eq:1taua_1e6_nua}
\end{equation}
Here the decoherence rate $1/\tau_a$ is proportional to frequency. Therefore we should expect shorter coherence time at higher frequencies. 
    
    
% \end{enumerate}

\section{Basics of nuclear magnetic resonance}\label{sec:intro:NMR}

After it was invented by Isidor Rabi in 1938, nuclear magnetic resonance (NMR) has developed to be one of the most precise techniques to measure magnetic field. NMR can work with zero- to ultra-low field (ZULF) which is well below earth field, and also field over \SI{10}{\tesla} where the magnet becomes the ceiling. It was natural to think of NMR as there is need to measure a pseudo-magnetic field precisely. 

\begin{figure}[ht]
    \includegraphics[width = \textwidth]{figures/01_Intro/TwoLevel_and_MagVec_20220929.png}
    \caption{Explanations on NMR mechanisms. (a) Zeeman splitting and two-level system. (b) Magnetization and precession because of excitation. }
    \label{fig:NMRmechanisms}
\end{figure}

\subsection{Zeeman splitting and polarization}\label{sec:intro:NMR:Zeeman_Pol}

For a nonzero spin in the magnetic field, Zeeman splitting lifts the degeneracy of in energy levels. In the microscopic picture, every 1/2 spin is in a superposition state of spin-up and spin-down, directions of which is defined as parallel or anti-parallel to the leading field. Depending on the magnetic moment, these two states correspond to ground and excited state in the language of atomic physics. Judged from its Hamiltonian in Eq.\,\eqref{eq:magfield_nucleonH}, the energy difference between two states is
\begin{equation}\label{eq:DeltaE}
    \Delta E = E_e - E_g =\hbar \gamma \mathrm{B} = 2 \mu \mathrm{B}\,,
\end{equation}
where $E_g$ and $E_e$ refer to the energy level for ground and excited state respectively, and $\mu$ is the magnetic moment. Therefore the angular frequency or so-called Larmor frequency $\omega$ is
\begin{equation}\label{eq:omega_gamma_B}
    \omega =2 \pi \nu = \gamma \mathrm{B} \,. 
\end{equation}
Sometimes we use frequency $\nu$ as well. 

Consider a typical proton NMR experiment where the temperature of sample is over \SI{100}{\kelvin} (\SI{-173.15}{\celsius}) and leading field is less than \SI{50}{\tesla}. Omitting the factor of $3/2$, kinetic energy $kT \geq \SI{8.61}{\milli\electronvolt} $, where $k$ stands for Boltzmann constant and $T$ stands for the temperature. The energy gap between two levels of NMR is $\hbar \gamma \mathrm{B} \leq \SI{8.8}{\micro\electronvolt}$. From the comparison of these two energies we can conclude that, unless for NMR at extreme conditions (low temperature or high field), the energy gap for NMR is well below the kinetic energy of particles:
\begin{align}
    \Delta E \ll kT\,.\label{eq:DeltaE_kT}
\end{align}
When at thermal equilibrium, the population of particles on either level obeys Boltzmann distribution. 
The ratio of the population at higher level over that at lower level can be given as
\begin{align}
    \exp\left(-\dfrac{E_e}{kT}\right) / \exp\left(-\dfrac{E_g}{kT}\right) = \exp\left(\dfrac{-\Delta E}{kT}\right)\,.
\end{align}
% Here we made the approximation based on Eq.\,\eqref{eq:DeltaE_kT}. 

In a NMR experiment with spin-1/2 sample, every spin could be flipped to the other state. Therefore, the difference between two populations matters. We define the ratio of this difference over whole population as polarization. When at thermal equilibrium, polarization $p$ can be computed as:
\begin{equation}\label{eq:thermalp}
	p =\left| \dfrac{N_{\uparrow} - N_{\downarrow}}{N_{\uparrow} + N_{\downarrow}} \right|
	= \left|\dfrac{\exp\left(-\frac{E_g}{kT}\right) - \exp\left(-\frac{E_e}{kT}\right)}{\exp\left(-\frac{E_g}{kT}\right) +\exp\left(-\frac{E_e}{kT}\right)}\right|
	=
	\dfrac{1 - \exp\left(-\frac{\Delta E}{kT}\right)}{1 +\exp\left(-\frac{\Delta E}{kT}\right)}\,. 
\end{equation}
Furthermore, we can simplify Eq.\,\eqref{eq:thermalp} based on the condition from Eq.\,\eqref{eq:DeltaE_kT}: 
\begin{align}
	p =\dfrac{\Delta E}{2kT} =\dfrac{\hbar \gamma \mathrm{B}}{2kT} = \dfrac{\mu \mathrm{B}}{kT}\,,\label{eq:thermalp_approx}
\end{align}
where  
\begin{align}
    \mu= \dfrac{\gamma }{2 \hbar}\,.
\end{align}


As we can tell from the Eq.\,\eqref{eq:thermalp_approx}, the polarization goes up with larger magnetic field. Nevertheless, even with the ectrme conditions mentioned earlier in this section, the polarization can only reach \SI{5e-4}{}. Typically the polarization can only reach $10^{-7}$ to $10^{-5}$. It is possible to boost polarization through other means. Spin exchange is one feasible method where the state of pre-polarized nucleus (e.g., proton of para-hydrogen, \ch{^{87}Rb} nucleus) is transferred to target nucleus. Polarization could be greater than 10 \%, which is approximately $10^6$ times larger than thermal polarization. 

\subsection{Magnetization and spin ensemble}

From here we start to emphasize the macroscopic picture of spin ensemble instead of individual ones. In other words, we focus on the collective effect of $\sim 10^{23}$ spins. The magnetization, or density of nuclear moment $\mathbf{M}$ is proportional to the magnetic field $\mathbf{B}$:
\begin{equation}
    \mathbf{M} = \chi \mathbf{B}\,.
\end{equation}
Here $\chi$ is the nuclear paramagnetic susceptibility, which is given by the Curie formula:
\begin{equation}
    \chi = \dfrac{j + 1}{3j} \dfrac{n \mu^2}{kT}\,,
\end{equation}
where $j$ is the spin number and $n$ is the spin density. Certainly we have $n = N_{all}/V$, where $N_{all}$ is the number of all spins, and $V$ is the volume of the NMR sample. Resultant magnetization from spin $j=1/2$ is
\begin{equation}
    \mathrm{M}_0 = \dfrac{n \mu^2}{kT} \mathrm{B} = \mu p N_{all}/V
    \,.\label{eq:M0_allspins}
\end{equation}
Here $p$ is taken from Eq.\,\eqref{eq:thermalp_approx}. Eq.\,\eqref{eq:M0_allspins} implies that, aside from thinking of nuclear paramagnetic susceptibility, the magnetization can also be regarded as originated from polarization. Notice that $\mathbf{M}$ is different from the concept of spin, which can be orientated towards any spatial direction. 

The interaction between magnetization and magnetic field can be described by Bloch equation\cite{Bloch_NuclearInduction_PhysRev.70.460}:
\begin{equation}\label{eq:BlochEq}
    \dfrac{d \mathbf{M}}{dt} = \gamma \mathbf{M} \times \mathbf{B}\,.
\end{equation}
The cross product indicates that when the orientation of $\mathbf{M}$ is not aligned to the leading field, it tends to precess around it. The Larmor frequency $\omega$ can also be drawn from here besides the theory of Zeeman splitting. 

\subsection{Spin-projection noise}\label{sec:intro:NMR:SPN}

By far we have discussed the polarization and magnetization in magnetic field, which are proportional to the number of total spins and magnitude of the field. However, early in 1946, Felix Bloch mentioned that due to statistically incomplete cancellation, $N$ nuclei can generate moment of the order $\sqrt{N}\mu$, even in the absence of orientation by external magnetic field\cite{Bloch_NuclearInduction_PhysRev.70.460}. The concept of spin-projection noise (SPN) was invented to describe this phenomenon. We explain the origin of it in two different ways in the following paragraph. 

To begin with, we recall the knowledge of magnetization properties. When at thermal equilibrium, the magnetization is along the leading field, and the expectation value of it on the transverse direction is zero. Nevertheless, the magnetization is not absolute steady. It fluctuates, resulting in a nonzero deviation from mean value. 

The other explanation involves the language of spin. Suppose we are measuring orientations of a spin ensemble along $\mathbf{e}_y$, and no magnetic field exists in this direction here. If they are spin-1/2 particles, the probabilities of obtaining projection on $+\mathbf{e}_y$ and $-\mathbf{e}_y$ are equal. But in measurements, the distribution of two results is not absolutely uniform, resulting in nonzero magnetization. 

Such deviation is described by the average absolute deviation (AAD). For normal distribution, the expectation value of AAD is $\sigma \sqrt{2/\pi}$, where $\sigma$ is the standard deviation. 
Following the thinking of spin-1/2 measurement, the results of measurements on single spins obey binomial distribution, where we have $\sigma = \sqrt{N}/2$ considering chances for both results are equal. For an ensemble of $N \approx 10^{23}$ spins, the results of measurements on net magnetization obey normal distribution since we have $1/\sqrt{N} \ll 1$. The standard deviation of this normal distribution is expected to be proportional to $\sqrt{N}/2$ as well. Therefore, the AAD of magnetization without existence of external magnetic field can be given by:
\begin{align}
    M = \dfrac{\sqrt{N}}{2} \sqrt{\dfrac{2}{\pi}} \dfrac{\mu}{V} = \sqrt{\dfrac{N}{2 \pi}} \dfrac{\mu}{V}\,.\label{eq:SPN_M}
\end{align}
Here $\sqrt{N/2 \pi}$ can be regarded as the number of effective spins (denoted as $N_{eff}$). For an ensemble of $N = 10^{22}$ spins, $N_{eff}$ induced by SPN is $\sqrt{N/2 \pi}=\SI{4e11}{}$. In contrast, thermal polarization of $p=10^{-7}$ generates $p N = 10^{15}$ effective spins, for which reason we are not influenced by SPN in typical NMR experiments. However, SPN could be significant when signal amplitude and system noise level are relatively low. 

\subsection{Relaxation}

The evolution of magnetization is not only governed by precession but also relaxation. It can be classified into longitudinal (along the leading field) and transverse (perpendicular to the leading field) relaxation. Usually z-direction is chosen as the orientation of leading field, thus the longitudinal component of magnetization is denoted as $M_z$, while the transverse one is $M_{xy}$. 

Relaxation drives magnetization towards equilibrium. During this process, $M_z$ evolves towards the magnitude that scales with thermal polarization, while $M_{xy}$ decays till zero. These two processes are described by relaxation time $T_1$ and $T_2^*$. 


\subsubsection{$T_1$ relaxation}

We can try to understand $T_1$ relaxation in such a process: Imagine that a nuclear spin-1/2 ensemble is moved into a magnetic field. %Before they are moved, the measured populations of spin-up and spin-down defined by the direction of leading field should be equal. 
Now probabilities of them appearing at each state are no more equal. The process of building up population difference (also, magnetization) requires the disturbance from sufficient molecule kinetic energy, interaction and time. 
After we take into account relaxation time, Bloch equation for $M_z$ becomes
\begin{align}
    % \dfrac{d M_x}{dt} &= \gamma (M_y B_z - M_z B_y) - \dfrac{M_x}{T_2^*}\label{eq:BlochMx}\\
    % \dfrac{d M_y}{dt} &= \gamma (M_z B_x - M_x B_z) - \dfrac{M_y}{T_2^*}\label{eq:BlochMy}\\
    \dfrac{d M_z}{dt} &= \gamma (M_x B_y - M_y B_x) - \dfrac{M_z - M_{0}}{T_1}\label{eq:BlochMz}\,,
\end{align}
where $M_{0}$ is the magnitude of magnetization at thermal equilibrium state. 
% We can look at the relaxation respectively. 
To better illustrate the relaxation, we assume the field to be simply $\mathbf{B} = \mathrm{B}_0 \mathbf{e}_z$. With the absence of magnetic field in $x$ and $y$ direction, the Bloch equation for $M_z$ is
\begin{align}
    \dfrac{d M_z}{dt} &= \dfrac{M_0 - M_z}{T_1}\,.
\end{align}
If the initial value of $M_z$ is zero, we can obtain the expression for this evolution as 
\begin{equation}
    M_z(t) = M_0(1 - e^{- t/T_1})\,. \label{eq:T_1relaxation_zeroMz0}
\end{equation}
From this expression we can officially define $T_1$ as the time for building up magnetization from zero to $M_0(1-e^{-1})$ through relaxation. Paramagnetic ions in solutions (like \ch{Cu+}) can shorten $T_1$ time in that they accelerate the exchange of atoms among molecules and thus equilibrium state can be established faster. 

Moreover, if the initial value of $M_z$ is an arbitrary value $M_{z0}$, we can obtain
\begin{equation}
    M_z(t) = M_0 + (M_{z0} - M_0)e^{- t/T_1}\,. \label{eq:T_1relaxation_NonzeroMz0}
\end{equation}
The plots for Eqs.\,\eqref{eq:T_1relaxation_zeroMz0} and \eqref{eq:T_1relaxation_NonzeroMz0} are included in Fig.\,\ref{fig:NMR_signal_configuration}(a). 

\begin{figure}[htp]
    \centering
    \includegraphics[width =\textwidth]{figures/01_Intro/M_and_T1_T2.png}
    \caption{NMR signal configurations. (a) (b) $T_1$ and $T_2$ processes with different initial values.}
    \label{fig:NMR_signal_configuration}
\end{figure}

\subsubsection{$T_2^*$ relaxation}

As mentioned above, the magnetization is the characterization of spin ensemble. If we look at smaller packets of the sample (still at macroscopic size), the individual packets experience different magnetic field. Firstly, the leading field $\mathbf{B}$ cannot be perfectly homogeneous throughout the volume. In addition, the spin-spin interaction influences homogeneity at microscopic scale. The differences between precession phases of the small packets grow over time. Such process is also referred to as ``dephase''. The macroscopic effect is that the net transverse magnetization decays. 
$T_2^*$ time is the product of these two mechanisms. The relaxation rate can be written as
\begin{equation}
    \dfrac{1}{T_2^*} = \dfrac{1}{T_2} + \dfrac{1}{T_{\delta}}\,,\label{eq:T2star_T2_Tdelta}
\end{equation}
where $T_2$ time takes into account the spin-spin interaction, and $T_{\delta}$ indicates inhomogenous broadening due to the variation of external magnetic field over the sample. In most of the experiments we do not have to differentiate them two except for spin echo experiment. But we have to keep this point in mind to avoid false expectation on $T_2^*$. One common mistake is to assume $T_2^*$ can be super long when we have perfect magnet, neglecting $T_2$ as the other bottleneck of $T_2^*$. Besides, due to $T_1$ relaxation, $T_2^*$ time is no greater than $T_1$. 

We know from Bloch equations that leading field $\mathbf{B} = \mathrm{B}_0 \mathbf{e}_z$ leads to Larmor precession. If we sit on the transverse magnetization vector, and spin with it at Larmor frequency, we do not observe the precession. In this so-called ``rotating frame'', we get rid of the leading field, and the Bloch equation for transverse magnetization (no magnetic field along transverse directions) is simplified to:
\begin{align}
    \dfrac{d M_{xy}}{dt} &= - \dfrac{ M_{xy} }{T_2^*}\,.
\end{align}
Therefore the expression for the evolution is
\begin{equation}\label{eq:T_2relaxation}
    M_{xy}(t) = M_{xy}(0) e^{- t/T_2^*}\,,
\end{equation}
with the initial value $M_{xy}(0)$ (see Fig\,\ref{fig:NMR_signal_configuration}(b) for the illustration). The expression implies exponential decay of the transverse magnetization due to relaxation. 
% Since the signal frequency is proportional to the field strength, the inhomogeneity of field reflects on the linewidth of NMR signal. NMR linewidth (full width at half maximum or FWHM) in the unit of \SI{}{\Hz} is inversely proportional to $T_2^*$:



\subsection{Chemical shift}

Due to the chemical environment and shielding from electron distribution, the nuclear spins could feel a different magnetic field:
\begin{equation}\label{eq:chemicalshiftfield}
    \mathrm{B}_{eff} = \mathrm{B} (1-\delta_c)\,,
\end{equation}
where effective field strength is different from external field $B$ by a factor of $\delta_c$. Consequently, the signal frequency is
\begin{equation}
    \omega = 2 \pi \nu = \gamma \mathrm{B}(1-\delta_c)\,. \label{eq:omega_gamma_B_chemicalshift}
\end{equation}
The common practice in the experiment is that we choose a value as reference frequency $\nu_{ref}$. The chemical shift can be computed as
\begin{align}
    \delta_c = \dfrac{\nu - \nu_{ref}}{\nu_{ref}}\,.
\end{align}
Since $\nu_{ref}$ is usually chosen close to the signal frequency $\nu$, thus we have $|\nu - \nu_{ref}| \ll \nu$ and $\delta_c = (\nu - \nu_{ref})/\nu$. This shift is usually at the order of $10^{-6}$. Therefore the unit of parts per million (\SI{}{\ppm}) is widely used here. 

\subsection{Spin-spin coupling}

When multiple non-zero nuclear spins exist in one sample molecule, the coupling between them influences the transition energy and resonant frequency. This is the so-called spin-spin coupling or J coupling. We have to account for additional terms in system Hamiltonian like:
\begin{equation}
    H_{l,k}  = J_{l,k} \mathbf{\sigma}_l \cdot \mathbf{\sigma}_k\,,
\end{equation}
where $J_{l,k}$ represents the coupling strength between spin $l$ and $k$. Depending on the orientations of spins, this term changes the transition energy for the system by $\pm E_{l.k}$, and signal frequency is shifted by $\pm E_{l.k}/h$. Therefore the splitting is a constant value for certain J coupling instead of relative separation scaling with external field. 

\subsection{Pulsed-NMR}\label{sec:intro:NMR:Pulsed-NMR}

% \subsubsection{Mechanism}

The idea of pulsed-NMR was first invented by Felix Bloch in 1946\cite{Bloch_NuclearInduction_PhysRev.70.460}. Here we look into the experimental scheme and the physical process, besides the spectroscopic analysis on NMR signal. 

\subsubsection{Experimental scheme}
Initially the leading field $\mathbf{B}_0$ polarizes sample and builds up magnetization $\mathbf{M}$. 

An excitation coil is placed next to the sample to shot pulse at resonant frequency. Such a pulse generates transient magnetic field $\mathbf{B}_1 = \mathrm{B}_1 \cos \omega t$ perpendicular to $\mathbf{B}_0$. Hence the magnetization is tilt by an angle from the direction of $\mathbf{B}_0$. 
In the rotating frame, $\mathbf{B}_1$ can be decomposed into two circularly polarized magnetic field: $\mathrm{B}_1 \cos (2\omega t)/2$ and $\mathrm{B}_1 \cos (0 t)/2$. The first one is far from resonance and not able to synchronize with $\mathbf{M}$, the effect of which is averaged to zero. 
The latter one acts as an DC field to $\mathbf{M}$. The tilt angle $\theta$ is proportional to the field strength and pulse duration: 
\begin{equation}\label{eq:tippingangle}
    \theta = \dfrac{1}{2} \gamma \mathrm{B}_1 t = \Omega_{R} t\,. 
\end{equation}
Here $\Omega_{R} = \gamma \mathrm{B}_1/2$ is the Rabi frequency. 
The transverse magnetization becomes: 
\begin{equation}\label{eq:MxyM0theta}
    M_{xy} = M_0 \sin \theta\,. 
\end{equation}
The optimal tilt angle for large signal amplitude is $\pi/2$. We call it $\pi/2$ pulse when magnitude of $\mathrm{B}_1$ and pulse duration can induce a tilt angle of \SI{90}{\degree}. Similarly we defined $\pi$ pulse to indicate a \SI{180}{\degree} tilt angle. In order to minimize the relaxation, the duration should be chosen to be small in comparison to $T_2^*$ and $T_1$ time. 

The pick-up coil oriented perpendicular to $\mathbf{B}_0$ can receive the precession signal from $M_{xy}$:
\begin{equation}
   M_{xy}(t) =  M_{xy0} e^{-t/T_2^*} \cos (\omega t + \varphi_0)\,,\label{eq:M_FID}
\end{equation}
where $M_{xy0}$ is the oscillation amplitude, $\varphi_0$ is the initial phase depending on the placement of coils, pulse strength and duration. This combination of precession and relaxation is called free induction decay (FID). The scheme of pulsed-NMR is illustrated in Fig.\,\ref{fig:Pulsed-NMR_scheme}. 

\begin{figure}[htbp]
    \centering
    \includegraphics[width = \textwidth]{figures/01_Intro/PulsedNMR_Schematic_20220911.png}
    \caption{Pulsed-NMR schematic. }
    \label{fig:Pulsed-NMR_scheme}
\end{figure}

\subsubsection{Spectroscopy analysis on pulsed-NMR signal}

Fourier transform (FT) is widely applied in analyzing oscillation signal. It transforms time domain to frequency domain, which enables us to distinguish signal of different frequencies:
\begin{align}
    \Tilde{s}(\nu) = \int_{-\infty}^{\infty} s(t) e^{-i2\pi \nu t} dt\,,
\end{align}
where $s(t)$ is the time-domain signal. 

In practice, we apply discrete Fourier transform (DFT) to process discrete data points. There are multiple ways to display the result of FT. 
Here we introduce power spectral density (PSD, also called power spectrum) $|\Tilde{s}(\nu)|^2$. It shows how signal average power is distributed in frequency domain. It is normalized to satisfy Parseval’s theorem: 
\begin{align}
    P = \dfrac{1}{T} \int_{0}^{T} |s(t)|^2 dt = \int_{0}^{\infty} |\Tilde{s}(\nu)|^2 d \nu\,,\label{eq:Parseval}
\end{align}
where $P$ is the averaged power and $T$ is measurement time. 
We can compute the average power of FID based on Eq.\,\eqref{eq:M_FID}:
\begin{align}
    P_{FID} = \dfrac{1}{T} \int_{0}^{T} M_0^2 e^{-2t/T_2^*} \cos^2 (2 \pi \nu t + \varphi_0) dt\,. 
\end{align}
If we have measurement time much longer than relaxation time ($T \gg T_2^*$), and relaxation rate much smaller than resonant frequency ($1/T_2^* \ll \omega/2\pi$), we can approximate calculation of the average power to:
\begin{align}
    P_{FID} &\approx \dfrac{1}{T} \int_{0}^{T} M_{xy0}^2 e^{-2t/T_2^*} dt \int_{0}^{T} \cos^2 (2 \pi \nu t + \varphi_0) dt \\
     &\approx  \dfrac{M_{xy0}^2 T_2^*}{4 T}\,, \label{eq:PFID}
\end{align}
from which we can learn that the average power of FID signal is proportional to $T_2^*$. The conditions for this equation hold true for CASPEr-gradient experiments. 

For pulsed-NMR FID signal, the PSD can be approximated by a Lorentzian function:
\begin{equation}
   Lorz(\nu) =  \dfrac{A \Gamma}{2\pi}\dfrac{1}{(\nu-\nu_s)^2 + \left(\Gamma/2\right)^2 }\,.
\end{equation}
Here $A$ is the area under the function curve, $\Gamma$ is linewdith (also FWHM) and $\nu_c$ is the central frequency. 
The amplitude of Lorentzian function at central frequency is 
\begin{equation}
    Lorz(\nu_c) =  \dfrac{2 A}{\pi \Gamma}\,,\label{eq:Lorz_nuc}
\end{equation}
which is proportional to area and the inverse of linewidth. 
For PSD of FID, the area is average power ($A = P$), which makes it convenient in signal power analysis. The central frequency is the resonant frequency $\nu_c = \gamma B / 2 \pi$. Besides the linewidth is determined by $T_2^*$:
\begin{equation}
    \Gamma = \dfrac{1}{\pi T_2^*}\,. \label{eq:Gamma_T2star}
\end{equation}
The normalized linewidth, or so-called fractional, is also widely accepted for characterizing NMR linewidth, definition of which is given by:
\begin{equation}
    \delta_s \equiv \dfrac{\Gamma}{\nu} = \dfrac{1}{\pi T_2^* \nu}
    \,.\label{eq:delta_s}
\end{equation}
Similar to chemical shift, we often use \SI{}{\ppm} as the unit for normalized linewidth. 

For FID signal larger $T_2^*$ means larger average power concentrated in a smaller frequency bin. Therefore we can expect signal amplitude proportional to $(T_2^*)^2$, which can be inferred from Eqs.\,\eqref{eq:PFID}, \eqref{eq:Lorz_nuc} and \eqref{eq:Gamma_T2star}:
\begin{align}
    |\Tilde{s}(\nu = \gamma B/2\pi)|^2 =  \dfrac{(M_{xy0} T_2^*)^2}{2 T} \,.
\end{align}

Two examples of \SI{2}{\Hz} FID power spectrum are shown in Fig.\,\ref{fig:PulsedNMR_Spectroscopy}. The chunk size for PSD is \SI{20}{\second}. The geometry of PSD follows what we computed in this section. 
\begin{figure}[htbp]
    \centering
    \includegraphics[width=\textwidth]{figures/01_Intro/PulsedNMR_Spectroscopy_20220911.png}
  \caption{Left: FID signal of $T_2^* = \SI{4}{\second}$ and $\SI{1}{\second}$. Right: PSD of FID signals. Black arrows indicate the linewidths. }
  \label{fig:PulsedNMR_Spectroscopy}
\end{figure}

Due to the presence of chemical shift, spin-spin coupling and geometry of leading field, the signal peak is not likely to be not exact Lorentzian. Structure of the signal peaks can reveal more information on the sample molecular composition and magnetic field configuration. 


\subsection{Spin echo}


Due to dephase, the signal amplitude of FID decays over time. We can rewrite FID signal in Eq.\,\eqref{eq:M_FID} with Eq.\,\eqref{eq:T2star_T2_Tdelta} into:
\begin{equation}
    M_{xy}(t) =  M_{xy0} e^{-t/T_2} e^{-t/T_{\delta}} \cos (\omega t + \varphi_0)\,. \label{eq:M_FID_T2_Tdelta}
\end{equation}
Spin echo aims at getting rid of the inhomogeneous broadening ($T_{\delta}$ relaxation) by pulse sequence. Here we discuss the experimental scheme along side method for $T_2$ measurement based on it. 

\subsubsection{Experimental scheme}

The key of spin ``echo'' is the pulses to refocus the magnetization during dephase. Here we again divide sample into smaller macroscopic packets, and denote the Larmor frequencies of $n$-th packet as $\omega_n$. At first, a $\pi/2$ pulse is applied to generate transverse magnetization, and the phase of $n$-th packet is a constant
\begin{align}
    \varphi_n = \varphi_{0}\,.
\end{align}
Then, the dephase happens as described in pulsed-NMR scheme. After time $\Delta t$ it becomes
\begin{align}
    \varphi_n^{\prime} = \omega_n \Delta t + \varphi_{0}\,.
\end{align}
At this moment a $\pi$ pulse is shot to flip the magnetization. The duration of the pulse is relatively short compared to $\Delta t$, and phase is changed into
\begin{align}
    \varphi_n^{\prime \prime} = k \pi - (\omega_n \Delta t + \varphi_{0})\,.
\end{align}
Depending on the orientation of the pulse, $k$ can be $0$ or $1$. 
After another time $\Delta t$, the phases once again turn into one constant:
\begin{align}
    \varphi_n^{\prime \prime \prime} = k \pi - (\omega_n \Delta t + \varphi_{0}) + \omega_n \Delta t = k \pi - \varphi_{0} \,.
\end{align}
In this way we refocus the phases after relaxation and hear the echo of spins. Another phrase to describe the process during the second $\Delta t$ is ``time reverse''. 

However, we neglected a few effects in the analysis. 
The first one is that $\omega_n$ is a not a constant. The motion of molecule make it experience the inhomogeneity of the leading field. The more important problem is the spin-spin interaction leading to time-dependent microscopic magnetic field, which is too chaotic to be compensated by pulse sequence. Therefore, we still have to live with $T_2$ relaxation, and the signal during $[\Delta t,\,3\Delta t]$ can be written in such a function based on Eq.\,\eqref{eq:M_FID_T2_Tdelta}:
\begin{align}
    M_{xy}(t) =  M_{xy0} e^{-t/T_2} e^{-|t -2\Delta t|/T_{\delta}} \cos (\omega t + \varphi_1)
    \,.
\end{align}
where $\varphi_1$ is an arbitrary phase. The scheme is illustrated in Fig.\,\ref{fig:SpinEcho_schematic}. 
% \begin{align}
%     M_{xy}(t) = \left\{
%     \begin{array}{cl}
%     M_{xy0} e^{-t/T_2} e^{-(2\Delta t - t)/T_{\delta}} \cos (\omega t + \varphi_1)
%     &  t \in [\Delta t, 2\Delta t)\\
%      & \\
%     M_{xy0} e^{-t/T_2} e^{-(t- 2\Delta t)/T_{\delta}} \cos (\omega t + \varphi_1)   
%     &  t \in (2\Delta t, 3\Delta t]\\
%     \end{array} \right.
%     \,,\label{eq:Spinecho_Dt3Dt}
% \end{align}


\begin{figure}[htbp]
    \centering
    \includegraphics[width=\textwidth]{figures/01_Intro/SpinEcho_Schematic_20220918.png}
    \caption{Spin echo experiment schematic. Blue stands for pulse and effects of them. Black circular arrows indicate the precession. Other colorful arrows stand for magnetization of small sample packets precessing at different Larmor frequencies. The ``echo'' or ``time reverse'' is indicated by $e^{-|t- t_c|}$. }
    \label{fig:SpinEcho_schematic}
\end{figure}

\subsubsection{$T_2$ measurement}

The echo amplitude decay with $e^{-t/T_2}$, for which reason we can measure $T_2$ with spin echo. The pulse sequence consists of one $\pi/2$ pulse at first and a series of $\pi$ pulses to refocus the magnetization. Between every two $\pi$ pulses, we can acquire precession signal. The signal sequence is also called train of echos as illustrated in Fig.\,\ref{fig:SpinEcho_schematic2}. Therefore we are able to do fitting to echo amplitude to obtain $T_2$. 
\begin{figure}[htbp]
    \centering
    \includegraphics[width=\textwidth]{figures/01_Intro/SpinEcho_Schematic2_20220918.png}
    \caption{Pulse sequence and train of echos.}
    \label{fig:SpinEcho_schematic2}
\end{figure}

\section{Experiment design for axion research}

\subsection{Measurement target}\label{sec:Experiment_Design:Measurement_target}

% There are two possible results for detection: we either discover axion-nucleon coupling, or exclude the coupling strength to some extend. 
As was discussed in \S\ref{sec:Axion_wind}, the axion detection is basically transformed into magnetic field measurement. If we choose the laboratory as $\mathbf{r} = 0$ for Eq.\,\eqref{eq:Ba_gaNN_rhoDM}, we can obtain the target field to measure:
\begin{align}
    \mathbf{B}_{a}(t) &\approx \mathrm{g_{aNN}} \sqrt{2 \hbar c \rho_{DM}} \sin(\omega_a t -\mathbf{k}\cdot\mathbf{r} + \phi) \mathbf{v}_a / \gamma \label{eq:Ba_vector}\,.
\end{align}

Here we first estimate the magnitude of $\omega_a$ and $\mathbf{B}_{a}$, and determine schemes for measurement afterwards. On one hand, we do not have strong limits on mass $m_a$ and corresponding frequency $\omega_a$ for axion. As a result, the frequency band to scan is wide open. On the other hand, we can take nuclear gyromagnetic ratios to be $\sim$ \SI{10}{\MHz/ \tesla} to estimate the magnitude of $\mathrm{B}_{a}$:
\begin{equation}
    \mathrm{B}_{a} \approx 10^{-7} \mathrm{g_{aNN}}\,. 
\end{equation}
Here $\mathrm{B}_{a}$ is in the unit of \SI{}{\tesla} and $\mathrm{g_{aNN}}$ is in \SI{}{\GeV^{-1}}. Considering that earth magnetic field is $\sim 10^{-5}\,\SI{}{\tesla}$, we expect the coupling $\mathrm{g_{aNN}}$ to be weaker than $10^{2}\,\SI{}{\GeV^{-1}}$. Otherwise the axion pseudo-magnetic field should be significant in earth field measurement. 

The velocity $\mathbf{v}_a $ is associated with the velocities of dark matter, Galaxy, solar system and Earth. Therefore the amplitude of axion signal ought to be modulated in the periods of day and year\cite{Gramolin_PhysRevD.105.035029}. 
Depending on the direction, the effect of axion field could be regarded as
\begin{enumerate}
    \item an offset field along the leading field of NMR, or
    \item an excitation field on the transverse plane to induce precession, or
    \item a combination of the two above. 
\end{enumerate}
We can either perform regular NMR to measure a time-dependent leading field $\mathrm{B}_0 + \mathrm{B}_{a}(t)$ for the first case, or measure the transverse magnetization to detect the excitation field for the second one. The thrid one could fall into either cases. 

\textbf{Leading field measurement}

We can perform regular pulsed NMR experiment to measure the field. An oscillating offset on the leading field makes Larmor frequency time-dependent, and in turn broadens the NMR linewidth.


Two conditions should be satisfied to resolve this broadening. To begin with, the original NMR linewidth should be close to or smaller than the broadening due to axion field, which is the amplitude of axion field $\gamma \mathrm{B}_a$. In other words, we need the variation in linewidth due to axion is greater than original variation due to relaxation:
\begin{equation}
    \gamma \mathrm{B}_a \geq  \Gamma\,.\label{eq:Leadingfield_1}
\end{equation}
% be close to or smaller than $\gamma \mathrm{B}_a$. With 

Moreover, the axion Compton frequency $\nu_a$ should be small compared to the relaxation rate. 
\begin{equation}
    \nu_a \leq \Gamma\,.\label{eq:Leadingfield_2}
\end{equation}
% To resolve NMR signal in PSD, the chunk size of time-domain data $T_{PSD}$ is usually chosen to be larger than relaxation time $T_2^*$. 
If axion field oscillates much faster than the relaxation, we are likely to find the average offset on signal frequency in the chunk size, which is zero. 

From Eqs.\,\eqref{eq:Leadingfield_1} and \eqref{eq:Leadingfield_2} we have:
\begin{equation}
    \gamma \mathrm{B}_a \geq \nu_a\,,\label{eq:offset_condition}
\end{equation}
indicating that the magnitude of $\mathrm{B}_{a}$ or equivalently, the coupling strength $\mathrm{g}_{aNN}$ limits the frequency band we can search for axion. 
% In Eq.\,\eqref{eq:B_a_magnitude} we derived $\gamma \mathrm{B}_a \approx 10^2 \mathrm{g_{aNN}}$. Therefore the frequency band we can search for axion is strongly limited by Eq.\,\eqref{eq:offset_condition}. 


\textbf{Excitation field measurement}

The other method is called a resonant experiment. When the Larmor frequency $\omega = \gamma \mathrm{B}$ is tuned to be on-resonance with axion Compton frequency $\omega_a$, axion field can be compared to the excitation field in pulsed-NMR scheme. Transverse magnetization can be induced by axion. In this way, the limitation becomes the sensitivity of detector and the magnitude of leading field. However, the disadvantage of resonant detection is the need to scan a large frequency band. The experiment procedure of CASPEr-gradient was developed based on the idea of measuring excitation field, which is stated in the subsequent section. 

\subsection{Procedure for CASPEr-gradient measurement}

Pulsed-NMR diagnostic is applied to acquire information of polarization, relaxation time and so on. Depending on the stability and properties of system, the pulsed-NMR diagnostic could be performed before, during or after the detection, or even multiple times.

The axion measurement at one frequency is basically the scheme for pulsed-NMR without actively applying pulses. The magnetization of nuclear spins on resonance with axion Compton frequency $\omega_a$ can be tilt by axion field, generating precession signal in the detector. Note that data acquisition time $T$ could be very long compared to regular pulsed-NMR. Raw time-domain data is sliced into smaller chunks of $T_{PSD}$ for DFT. The power spectra obtained with DFT are averaged to decrease the noise level.
% The time-domain data is later sliced into smaller chunks for generating and averaging power spectrum. 
On one hand, in the case of discovering signal, more statistical analyses are to be done to verify the origin of it. 
On the other hand, if we do not find signal in the averaged power spectrum, we can exclude the coupling constant to certain extend at the Larmor frequency. 
Lastly, if the target frequency band is too large to be covered with one measurement, we are going to repeat it at multiple frequencies. The procedure is shown in Fig.\,\ref{fig:Procedure}. 

\begin{figure}[htp]
    \centering
    \includegraphics[width = \textwidth]{figures/01_Intro/CASPEr_Procedure_20220914.png}
    \caption{Procedure of CASPEr experiment. Step 1 to 5 are measurement at one single frequency. Step 6 indicates repetition of the steps before at different frequencies. The depths in spectra during step 5 and 6 are exclusion in mass-coupling parameter space. The lineshape depends on the axion linewidth, frequency-bin size and NMR linewidth. }
    \label{fig:Procedure}
\end{figure}





